\chapter{Against God}
There is a recurring myth in the ancient Mediterranean. In it, the Shining One (Hebrew: Helel, Greek: Phaethon) tries to usurp the Sun or the supreme deity, and is cast down or punished for his presumption. This is a familiar myth in our culture, due mostly to the Greek version. The Semitic version is less well known, in large part because one of the few written references we have to it has been lost in translation: 

"How you have fallen from heaven, morning star, son of the dawn! You have been cast down to the earth, you who once laid low the nations! You said in your heart, 'I will ascend to the heavens; I will raise my throne above the stars of God; I will sit enthroned on the mount of assembly, on the utmost heights of Mount Zaphon. I will ascend above the tops of the clouds; I will make myself like the Most High.' But you are brought down to the realm of the dead, to the depths of the pit." \cite{ref122}

The English term "morning star" is being used to translate the Hebrew Helel. We can imagine the mythology here fairly easily—the brightest star in the sky, refusing to share its place with the other stars, and instead jumping up into the sky at dawn, ahead of the sun. Then at sunrise it is wiped away, only for the story to repeat the next day, an endless cycle of celestial hubris.

Of course, most of us are more familiar with another translation, the King James, and another variant of the myth, which uses the Medieval Latin name for the morning star: Lucifer.

And that's it. That is the entirety of the Biblical story of Lucifer. Everything else is folklore and tradition, which is to say, fanfiction: that Lucifer was an angel, that he is the entity referred to as Satan in the book of Job, that he is the serpent in Genesis, that he is the Beast or the Dragon in Revelation; none of this is actually stated in the Bible itself, which just gives the story of a proud figure who rises up and is cast down. (There is nothing wrong with this, of course. Sacred texts are just one element in the complex of ideas, behaviors, and institutions that is a religion \cite{ref123}.) 

If we are to look for such a figure of hubris in Rebellion, Kyubey is pretty clearly that figure. He explicitly states that his goal in placing Homura within the barrier is to "interfere with"—that is, control and usurp—Madoka, who created the present universe, and as a result he is furiously punished by Madoka's herald Homura for his crimes. The result is a new universe in which, the stinger seems to imply, Kyubey's power has been almost completely stripped by the presence of Homura as an active and engaged demiurge.

There is another reading available, however, if we look at one of the most famous "fanfiction" versions of Lucifer, Paradise Lost. In discussing Milton's epic poem, however, it is important first to understand what an epic poem is. Understood cladistically, we can view the epic as a genre mostly descended from the works of Homer; the usual definition provides a list of common generic traits in terms of subject matter and structure, of which the most important for our purposes are that it involves events occurring on a national, cosmic, or global scale; follows the exploits of a larger-than-life, often supernaturally empowered hero; and utilizes a distinctive style that elevates it above normal discourse. In addition, epics usually start with an invocation and declaration of theme, begin in medias res, and contain lengthy monologues, often at least one flashing back to describe events prior to the opening. \cite{ref124}

Part of what makes Paradise Lost such a fascinating read is that Satan is consistently an incredibly vile character, a lying, cheating, self-serving manipulator—but he is also the epic hero upon whose adventures the story focuses. He is capable of being extremely charming and persuasive, to the point of convincing some very important critics (most famously Blake, who opined that Milton was "of the devil's party without knowing it" \cite{ref125}) that his cause is actually right. Keep in mind, said cause is the conquest of the world and enslavement and extermination of humanity! (All of which he succeeds at. Satan's son-grandson Death and daughter-bride Sin create a bridge from Hell to Earth at the end of the story, the Fall subjects humanity to Sin, and Adam and Eve are punished with mortality, which they pass on to their descendants, killing us all \cite{ref126}.) 

These contradictions have the effect of making Satan a morally ambiguous figure in a sense—he is structurally heroic but diegetically villainous, essentially. But what does this have to do with Rebellion?

Well, consider: Rebellion opens with an invocation of the Law of Cycles and a statement of theme, the cycle of despair in an irredeemable world and the escape into oblivion. We are then dropped into the middle of a situation that does not follow from the end of the series at all, forced to wait until some lengthy exposition by Kyubey much later in the movie to find out what happened to create this circumstance. That our heroine is supernaturally empowered goes almost without saying, as she is a magical girl—but she is empowered beyond that by her status as the witch in whose labyrinth the action of the series takes place, and then later by her love, which transforms her into a demon. In the process, she expands the story to be cosmic in scope. And as for a distinctive style that elevates it above the normal discourse of anime film, well, see Chapter 16.

So Rebellion is an epic. But more than that, it is the epic of how Homura went from being Madoka's "very best friend" in Episode 12 to calling her an enemy at the end of this movie. It is the epic, in other words, about how the closest and most loyal follower of the closest thing Madoka Magica has to a goddess fell to become a demon, and at the same time conquered the material universe with the stated intent of shepherding it to its destruction—yet throughout, remains a morally ambiguous figure, such that debates still rage across Internet fora as to her moral status.

Just as the series is not merely about a Faustian bargain, but actually in many ways retraces the story of Goethe's Faust, including some fairly obscure elements such as time travel, so too is Rebellion more than just a hubristic fall; structurally and in its juxtaposition of the epic hero with the moral fall, it resembles Paradise Lost.
